\chapter{The Problem}

\section{The operating system is the bottleneck}

Critical software has two requirements that pull in opposite directions: it must be
\textbf{correct} (failures are expensive or dangerous), and it must be
\textbf{performant} (latency and throughput matter). Modern systems try to satisfy
both by piling layers onto an operating system whose fundamental abstractions were
designed for timesharing, not for the workloads we run today.

The result is what we call the \textbf{retrofit tax}: the cost, in complexity and
performance, of making an OS designed for synchronous, trust-everyone, monolithic
execution behave as if it were asynchronous, least-privilege, and decomposed.

\section{The retrofit tax, itemized}

\subsection{Async over blocking}

Most operating systems expose synchronous system calls: \texttt{read()} blocks the
calling thread until data arrives. When applications need concurrency without
dedicating one thread per operation, they reach for \texttt{epoll}, \texttt{kqueue},
or \texttt{io\_uring}. These are effective, but they are \emph{retrofits}: the kernel
internally still blocks on disk I/O, network I/O, and locks, while the completion
mechanism translates those blocking events into readiness notifications for
userspace.

This translation layer brings with it:
\begin{itemize}
    \item Async-signal-safety problems --- the kernel may deliver a
        completion while userspace holds a lock, requiring elaborate
        signal-handler discipline or a dedicated thread pool.
    \item Thread-pool overhead --- many async runtimes spawn a pool of OS
        threads to convert blocking syscalls into futures, undoing the
        resource-efficiency argument for async in the first place.
    \item Two concurrency models --- the application uses async/await; the
        kernel uses threads and blocking. The mismatch forces every boundary
        crossing to bridge two foreign paradigms.
\end{itemize}

\subsection{Ambient authority}

A conventional OS process has access to the filesystem namespace, the network, and
other global resources by default. Security is subtractive: you start with
everything and try to take things away (chroot, seccomp, namespaces). This is
brittle because it requires the developer to enumerate what \emph{should not} be
accessible --- an unbounded problem. In practice, most programs run with far more
authority than they need, and exploits exploit that gap.

\subsection{Hardware asynchrony disguised as software synchrony}

Hardware is inherently asynchronous with respect to the CPU. A disk controller
signals completion via an interrupt. A network card delivers packets on its own
schedule. A timer fires after a programmed interval. Yet the operating system
interface presents these as synchronous calls: \texttt{read()} returns when the
data is ready, blocking the calling thread in the meantime.

This is a leaky abstraction. The thread that blocks on I/O is an OS resource that
cannot be used for anything else while it waits. High-concurrency servers
compensate with thread pools; real-time systems compensate with carefully tuned
polling loops. Neither is the natural solution.

\subsection{Fault propagation across trust boundaries}

A bug in a device driver can crash a monolithic kernel. A memory corruption in one
application can, in a system without hardware-enforced isolation, corrupt another.
Conventional microkernels solve this with address-space isolation, but at the cost
of IPC overhead that often drives designers back toward monolithic architectures
for performance.

The industry has settled on a compromise: run untrusted code in virtual machines;
run trusted code in the kernel; hope the trusted code has no bugs. For critical
systems, hope is a poor strategy.

\section{Why Rust ownership is not enough}

Rust's type system eliminates data races, use-after-free, and null-pointer
dereferences at compile time. This is a transformative improvement for systems
programming, but Rust programs still run on operating systems that do not share
Rust's ownership discipline.

\begin{itemize}
    \item A Rust program cannot protect itself from a C driver that
        corrupts kernel memory. Isolation must be enforced by the hardware
        (MMU, IOMMU), not by the language.
    \item A Rust \texttt{Mutex<Vec<T>>} crossing a syscall boundary is
        still a shared-memory lock with all the contention, priority inversion,
        and deadlock risks that implies. The kernel does not know the lock exists.
    \item Rust's \texttt{async}/\texttt{await} compiles to a state machine
        driven by a userspace executor, but the executor must still interact with
        the kernel's blocking I/O model. The \texttt{Future} trait and the
        kernel's wait queue speak different languages.
\end{itemize}

The insight is: Rust eliminates \emph{some} classes of bugs within a single
address space. But the OS defines the rules for \emph{between} address spaces ---
and if those rules are "any process can do anything until restricted," Rust's
guarantees are local optimizations on a global problem.

\section{What needs to change}

We need an operating system whose foundational abstractions align with the demands
of critical, performant software:

\begin{enumerate}
    \item Asynchronous by construction: Operations that may wait for
        external progress should have submission/completion semantics, not
        blocking-call semantics. No subsystem should force a thread to dedicate
        itself to one operation.

    \item Capability-based, not ambient, authority: A process should be
        able to perform an operation only if it holds a capability authorizing it.
        Capabilities are unforgeable and must be explicitly granted.

    \item Isolation by default: Every protection domain should run in its
        own address space. Communication crosses a kernel-mediated boundary.
        A fault in one domain should not propagate to another.

    \item Ownership that crosses boundaries: Data movement between
        domains should transfer ownership, not copy bytes. The kernel should move
        page mappings, so the sender loses access and the receiver gains it.

    \item Bounded resource consumption: Every queue, table, and buffer
        should be bounded. Backpressure should propagate explicitly. No unbounded
        allocation should be a correctness requirement.

    \item The OS and the language runtime should agree: The kernel's model
        of concurrency (asynchronous, completion-driven) should match the
        language's model of concurrency (futures, wakers). No translation layer
        should be needed.
\end{enumerate}

CharlotteOS is an operating system built from these six principles. This manual
explains the architecture, the programming model, and --- most importantly --- why
the combination eliminates the retrofit tax and provides a strong foundation for
software that must be both correct and fast.

\endinput
